\documentclass[11pt]{article}
\usepackage[T1]{fontenc}
\usepackage{lmodern,microtype,amsmath,amssymb,amsthm,booktabs,array}
\usepackage[margin=1in]{geometry}
\usepackage[hidelinks]{hyperref}
\hypersetup{pdftitle={A chain-cycle separation for leximin protection in three-sided assignment markets},pdfsubject={Anonymous manuscript for specialist review}}
\usepackage{enumitem}
\setlist{nosep}
\newtheorem{theorem}{Theorem}[section]
\newtheorem{lemma}[theorem]{Lemma}
\newtheorem{proposition}[theorem]{Proposition}
\newtheorem{corollary}[theorem]{Corollary}
\theoremstyle{remark}\newtheorem{remark}[theorem]{Remark}
\newcommand{\core}{\mathcal C}
\newcommand{\lex}{\operatorname{leximin}}
\newcommand{\R}{\mathbb R}
\newcommand{\N}{\mathbb N}
\title{A chain--cycle separation for leximin protection\\in three-sided assignment markets}
\author{}
\date{September 6, 2026}
\begin{document}
\maketitle
\begin{abstract}
We compare a player's leximin payoff with its largest payoff in the same
assignment-game core. In complete strict three-sided markets with nonnegative
pair values along a chain, every player receives at least one third of this
core ideal. The constant is sharp. Allowing the third pair interaction
eliminates every positive guarantee independent of market size. For each
integer $q\geq1$, we construct a market with $9q+3$ players, pair entries in
$\{0,1,2\}$, and a strictly positive core allocation, in which a player's
full leximin payoff is $1/[3(q+1)]$ while its core ideal is one. The chain
bound follows from a classical potential-slack balancing certificate whose
cycles contain three payoff arcs. The cyclic construction is certified by
two explicit full-core allocations and a disjoint partial matching one unit
below optimum. The proofs are valid for every parameter value; exact rational
checks and optional raw-core satisfiability checks accompany the manuscript.
\end{abstract}

\noindent\textbf{Keywords:} assignment game; core; leximin; ideal payoff;
pair-additive value; lexicographic slack balancing.

\section{Introduction}
A core allocation protects every coalition against a profitable deviation.
Leximin additionally favors the worst-paid individuals, successively through
all payoff levels. These two requirements do not by themselves say how much
of an individual's feasible bargaining range is retained. For a full core
$\core$ and its full ascending leximin allocation $x$, we study the fraction
$x_i/H_i$, where $H_i=\max\{y_i:y\in\core\}$ is the ideal payoff of player $i$.
The maximization defining $H_i$ permits all other payoffs to change.

Our comparison keeps the number of sectors fixed at three. Only complete
triples produce value; a pair does not earn its pair contribution on its own.
Each triple value is the sum of nonnegative values on its three pair
interactions. A chain uses two of these interactions, while the full class
also permits the third. We prove that the best uniform ideal-protection
fraction is $1/3$ for chains and zero for the full pair-additive class.
Thus a restriction on the interaction structure has a consequence that
cannot be recovered from nonemptiness of the core alone.

\paragraph{Contributions and established ingredients.}
The main construction realizes a vanishing fraction using a fixed integer
alphabet of pair values, three equal sectors, and a core allocation positive
in every labeled coordinate. Its target is already fixed on the first
optimal leximin face, so no arbitrary choice among max-min optimizers is
involved. The positive chain theorem provides the comparison: a
self-contained specialization of classical selected-arc slack balancing
gives the factor three. The underlying fairness benchmark, assignment
models, LP duality, and slack-balancing method are established ingredients.
We do not claim a new general fairness criterion or the first vanishing
example across all cooperative games. Historical priority of the precise
restricted-class separation is the question addressed by the accompanying
source-comparison and specialist-review packet; no external priority
endorsement is asserted here.

\paragraph{Related work.}
Multisided assignment cores have a substantial literature, including
Quint~\cite{Quint1991} and Lucas~\cite{Lucas1995}. The strict
supplier--firm--buyer model is associated with Stuart~\cite{Stuart1997};
Atay and Solymosi~\cite{AtaySolymosi2018}, Definition~3 and equations~(1)--(2),
give a precise later formulation allowing signed adjacent-pair values and
replacing negative triple totals by zero. Their core and bargaining-set
results concern stability sets, rather than the coordinate-ideal fraction
studied here. Nonnegative chain values are a subclass of that model.
Generalized assignment markets in which smaller teams also produce value
have additional coalition constraints and must be distinguished from our
strict model~\cite{AtayNunez2020}.

Ideal-based core protection is explicit in N\'u\~nez and
Rafels~\cite{NunezRafels2004}, Definition~6, in their bilateral bargaining
analysis. A unit-capacity path in Gangam et al.~\cite{Gangam2025},
Section~3.1.4, has uniform equitable payoff $1/r$ on its $r$ edges.
Its full core is the unit simplex, so each ideal is one: vanishing
protection in a broader game class follows immediately. That path uses a
productive coalition whose size grows with $r$. Our construction instead
uses only strict triples and pair values in $\{0,1,2\}$.

The potential argument for chains belongs to classical lexicographic
slack balancing. Held et al.~\cite{Held2003}, Section~3.2, explicitly
optimize ascending-sorted slacks on selected arcs while retaining other
potential inequalities. We use a cut perturbation and cycle telescoping
to give the needed optimality certificate directly. This is an
application-specific proof, not a methodological novelty claim.

Extreme-core descriptions for assortative multisided
markets~\cite{Martinez2020} and supermodular markets with local
monotonicity~\cite{Martinez2026} impose restrictions different from
nonnegative pair additivity. We make no assertion that arbitrary extreme-core
or transportation representations preserve the entire labeled payoff core,
the raw sorted selector, and our pair-value restrictions. The distinct
2023 manuscript on extreme core points~\cite{Martinez2023} is now available
for comparison. Its Theorem~2 redescribes the full core through cumulative
sums of auxiliary-game core allocations, retaining nonnegativity and the
triple inequalities outside a central strip. Theorems~3--4 simplify this
under additional restrictions. These are core descriptions; they do not
state the vanishing family or its full-leximin ideal ratio. The original
texts of~\cite{Quint1991,Lucas1995,Stuart1997} remain access gaps at the
date of this version. No theorem below relies on an uninspected statement
from these originals.

\section{Model and the separation theorem}
Let $I,J,K$ be disjoint sectors of equal finite size $m\geq1$, and put
$N=I\cup J\cup K$. Every triple $(i,j,k)\in I\times J\times K$ is
available and has nonnegative value $w_{ijk}$. A matching is a collection of
player-disjoint triples, possibly empty. For $S\subseteq N$, define
\[
 v(S)=\max\left\{\sum_{(i,j,k)\in M}w_{ijk}:
       M\text{ is a matching contained in }S\right\}.
\]
In particular, singletons and pairs have value zero. Write $T=v(N)>0$.
The full payoff core is
\begin{equation}\label{eq:core}
 \core=\left\{p\in\R_+^N:\sum_{a\in N}p_a=T,
                \quad p_i+p_j+p_k\geq w_{ijk}
                \ \forall(i,j,k)\right\}.
\end{equation}
This is exactly the coalitional core: necessity follows from individual
rationality and triple coalitions, and sufficiency follows by summing
triple inequalities over a matching of any coalition and using
nonnegativity for unmatched players. Equal sector sizes and nonnegative
values allow every partial matching to be completed without loss, so
$T$ is attained by a perfect matching.

For $p\in\R^N$, let $p^\uparrow$ be its coordinates sorted in
nondecreasing order. The full leximin allocation is the point $x\in\core$
that lexicographically maximizes $p^\uparrow$ over $p\in\core$.
Throughout, all labeled coordinates remain in this optimization, including
coordinates with zero ideal.

\begin{lemma}\label{lem:unique}
Every nonempty compact convex subset of $\R^N$ has a unique labeled full
leximin allocation.
\end{lemma}
\begin{proof}
Each sorted coordinate is continuous. Maximizing these coordinates
successively over nested nonempty compact sets gives existence.
For uniqueness, suppose distinct $x,y$ have the same optimal sorted vector.
Let $\alpha$ be the smallest value among $x_a,y_a$ at coordinates where
they differ. At every coordinate with a value below $\alpha$, the two
entries agree. At every remaining
coordinate, $(x_a+y_a)/2\geq\alpha$, and equality holds only when both
entries equal $\alpha$. Since $x$ and $y$ have the same multiplicities of
each value and differ at level $\alpha$, their average has strictly fewer
entries at that level, with all lower levels unchanged. Its sorted vector
is lexicographically larger, contradicting convexity and optimality.
\end{proof}

When $\core\ne\varnothing$, it is compact by~\eqref{eq:core}.
Set $H_a=\max_{p\in\core}p_a$ and only form $x_a/H_a$ when $H_a>0$.
A \emph{chain market} has
\begin{equation}\label{eq:chain}
 w_{ijk}=P_{ij}+Q_{jk},\qquad P,Q\geq0.
\end{equation}
A \emph{pair-additive market} has
\begin{equation}\label{eq:pair}
 w_{ijk}=A_{ij}+B_{ik}+C_{jk},\qquad A,B,C\geq0.
\end{equation}
The chain class is contained in the pair-additive class after renaming the
matrices. Let $c_{\rm chain}$ and $c_{\rm pair}$ be the infima of $x_a/H_a$
over the respective classes with $T>0$, nonempty full core, and $H_a>0$.

\begin{theorem}[Chain--cycle separation]\label{thm:main}
\hfill
\begin{enumerate}[label=(\roman*)]
\item Every chain market has nonempty core, and its full leximin allocation
      satisfies $x_a\geq H_a/3$ for every player. The constant is sharp.
\item For every integer $q\geq1$, there is a pair-additive market with
      $9q+3$ players, all pair entries in $\{0,1,2\}$, and every triple
      value at most $5$, with integral worth $T=6q+1$, a strictly positive
      core allocation, and a target player satisfying
      \[
             x_a=\frac{1}{3(q+1)},\qquad H_a=1.
      \]
      The target is fixed at this value throughout the entire first
      optimal max-min face of the core.
\end{enumerate}
Consequently $c_{\rm chain}=1/3$ and $c_{\rm pair}=0$.
\end{theorem}

Part~(ii) shows that \emph{allowing} a cyclic interaction removes the
class-wide guarantee. It does not assert that every cyclic market has weak
protection, or classify interaction graphs with more than three sectors.
The fixed integer alphabet is substantive: merely bounding real weights
could be achieved by rescaling each instance without changing $x_a/H_a$.

\section{The chain bound}
\subsection{The entire core as a potential system}
Choose maximum-weight perfect matchings in $P$ and $Q$. Their values sum
to $T$: any triple matching projects to two bilateral matchings, and the
two chosen matchings can be joined through their common middle players.
Relabel the outer sectors independently so both chosen matchings are
diagonal, and set $s_j=P_{jj}+Q_{jj}$. Write payoffs $(u,v,z)$ in sector
order. At every core point, summing the diagonal triple inequalities and
using efficiency makes each one tight. Thus
\begin{equation}\label{eq:diag}
                 u_j+v_j+z_j=s_j.
\end{equation}
The full core is precisely the set defined by
\begin{align}
 u_j,z_j&\geq0, & v_j=s_j-u_j-z_j&\geq0,\label{eq:nonneg}\\
 u_i-u_j&\geq P_{ij}-P_{jj} &&(i,j),\label{eq:U}\\
 z_k-z_j&\geq Q_{jk}-Q_{jj} &&(j,k).\label{eq:Z}
\end{align}
Indeed, the triples $(i,j,j)$ and $(j,j,k)$ imply the two differences;
adding them and substituting~\eqref{eq:diag} gives the inequality for an
arbitrary triple. Totals in~\eqref{eq:diag} add to $T$.

Nonemptiness follows also from standard bipartite matching duality. The
maximum-weight matching LP with incidence constraints at most one has
nonnegative dual vertex prices and an integral optimum. For nonnegative
weights on a complete equal-sided graph, that optimum is also the perfect
matching optimum. Take optimal nonnegative covers $(\alpha_i,\beta_j)$
for $P$ and $(\gamma_j,\delta_k)$ for $Q$. Then
$(\alpha_i,\beta_j+\gamma_j,\delta_k)$ is a nonnegative triple cover
of total $T$, hence a point of~\eqref{eq:core}. This uses classical
matching integrality and LP duality, not a new balance criterion.

\subsection{A cycle certificate for selected payoff slacks}
Consider a finite directed graph, a distinguished anchor with potential
zero, and slacks
\[
       \sigma_e(q)=c_e+q_{\mathrm{tail}(e)}-q_{\mathrm{head}(e)}\geq0.
\]
Some arcs are designated as player payoffs; the others are structural
constraints. Parallel arcs are allowed. We use the following elementary
optimality certificate for classical selected-slack balancing.

\begin{lemma}\label{lem:cycle}
Suppose a feasible potential vector realizes the full leximin payoff $x$.
For a payoff arc $e$ with distinct endpoints, there is a simple directed
cycle containing $e$ such that all its structural arcs are tight at $x$
and all its payoff arcs $f$ satisfy $x_f\leq x_e$. If that cycle has $r$
payoff arcs, then every feasible payoff vector $y$ satisfies $y_e\leq r x_e$.
\end{lemma}
\begin{proof}
Put $\alpha=x_e$. Form a directed subgraph $D$ consisting of the tight
structural arcs and all payoff arcs with value at most $\alpha$. Suppose
there is no path in $D$ from the head of $e$ to its tail. Let $S$ be the
nodes that can reach the tail in $D$, including the tail itself. The head
does not belong to $S$. No arc of $D$ enters $S$ from its complement.

Set a potential direction $h$ equal to one on $S$ and zero outside it;
subtract its anchor value everywhere to preserve the anchor. Every arc
of $D$ then has nonnegative slack derivative, and $e$ has derivative one.
All omitted structural arcs have strictly positive slack. All omitted
payoff arcs have value strictly above $\alpha$. Since there are finitely
many arcs, a sufficiently small positive step preserves feasibility and
keeps every decreasing payoff above $\alpha$. Let $k$ be the number of
payoffs initially at most $\alpha$. None of the first $k$ order statistics
can decrease: selected payoffs do not decrease, and omitted payoffs stay
above $\alpha$. The target moves above $\alpha$, leaving at most $k-1$
payoffs at or below that level. The $k$th order statistic therefore
strictly increases, a contradiction.

There is consequently a head-to-tail path in $D$. Take a simple one.
It cannot contain $e$, since that would revisit its starting head or
reach its terminal tail before the end. Appending $e$ gives a simple
cycle $Z$. Potential differences telescope around this cycle. At $x$
its structural slacks vanish, so for every feasible $y$,
\[
 \sum_{\substack{f\in Z\\f\text{ payoff}}}y_f
 \leq\sum_{f\in Z}c_f
 =\sum_{\substack{f\in Z\\f\text{ payoff}}}x_f
 \leq r x_e.
\]
The first inequality discards nonnegative structural slacks at $y$.
Nonnegativity of the other payoff slacks then gives $y_e\leq r x_e$.
\end{proof}

Neither connectedness of $D$ nor positive payoff levels were assumed.
If $x_e=0$, the same certificate implies that the coordinate is zero at
every feasible point.

\subsection{Exactly three payoff arcs}
Introduce nodes $0,U_j,Z_j$ with potentials
$q_0=0$, $q_{U_j}=u_j$, and $q_{Z_j}=-z_j$.
The payoff arcs are
\begin{center}
\begin{tabular}{lll}\toprule
Arc & Constant & Slack\\\midrule
$U_j\to0$ & $0$ & $u_j$\\
$0\to Z_j$ & $0$ & $z_j$\\
$Z_j\to U_j$ & $s_j$ & $v_j=s_j-u_j-z_j$\\\bottomrule
\end{tabular}
\end{center}
Structural arcs $U_i\to U_j$ have constants $P_{jj}-P_{ij}$;
structural arcs $Z_j\to Z_k$ have constants $Q_{jj}-Q_{jk}$.
Their nonnegative slacks give~\eqref{eq:U}--\eqref{eq:Z}.
Diagonal structural loops may be omitted. Structural constants may be
negative; only slacks are constrained to be nonnegative.

Every simple directed cycle containing a payoff arc has exactly three
payoff arcs. Structural arcs remain in one layer. The only transition
from the $Z$ layer to the $U$ layer is a middle-player arc; the return
passes through $0$, using one arc from each outer sector. A simple cycle
visits $0$ at most once, so it cannot make two such transitions.
Lemma~\ref{lem:cycle} gives $y_a\leq3x_a$ for every $y\in\core$.
Maximizing over $y$ proves part~(i) of Theorem~\ref{thm:main}.

For sharpness, use one player in each sector with $P_{11}=1$ and
$Q_{11}=0$. Its full core is the three-coordinate unit simplex,
$x=(1/3,1/3,1/3)$, and every ideal is one. Thus no larger uniform
fraction is possible.

\begin{corollary}\label{cor:distance}
For every chain market and every $y\in\core$,
\[
                    \frac{\|x-y\|_1}{T}\leq\frac43.
\]
This bound is attained for a finite market.
\end{corollary}
\begin{proof}
The theorem implies $\min(x_a,y_a)\geq y_a/3$ for every $a$.
Since both totals are $T$,
$\|x-y\|_1=2(T-\sum_a\min(x_a,y_a))\leq4T/3$.
The preceding simplex and $y=(1,0,0)$ give equality.
\end{proof}

\section{A cyclic family with vanishing protection}
Fix an integer $q\geq1$. Use the common team-label set
\[
           \mathcal T=\{L_0\}\cup\{O_s,L_s,R_s:1\leq s\leq q\}.
\]
For every label $t$, create players $a_t,b_t,c_t$ in the respective
sectors. Each sector has $3q+1$ players. Define
$w(a_t,b_u,c_v)=A_{tu}+B_{tv}+C_{uv}$ by the following nonzero entries;
all unspecified pair entries are zero.
\begin{center}
\begin{tabular}{lll}\toprule
Pair entries & Range & Value\\\midrule
$C_{L_s,L_s}$ & $0\leq s\leq q$ & $1$\\
$C_{R_s,R_s}$ & $1\leq s\leq q$ & $1$\\
$A_{O_s,O_s},\ B_{O_s,O_s}$ & $1\leq s\leq q$ & $2$\\
$A_{O_s,R_s},\ B_{O_s,L_s}$ & $1\leq s\leq q$ & $2$\\
$C_{R_s,L_s},\ C_{L_{s-1},O_s}$ & $1\leq s\leq q$ & $1$\\\bottomrule
\end{tabular}
\end{center}
There are $8q+1$ nonzero entries. The diagonal perfect matching consists
of $(a_t,b_t,c_t)$ for all $t\in\mathcal T$. Each $L$ or $R$ triple has
value one and each $O$ triple has value four, so this matching is worth
$6q+1$.

\subsection{Two explicit full-core covers}
Put $\tau=1/[3(q+1)]$. Define allocations $u$ and $y$ in $(a,b,c)$ order:
\begin{equation}\label{eq:witnesses}
\begin{array}{c|c|c}
\text{team}&u&y\\\hline
L_s&\bigl(\tau,\ 1-(3s+2)\tau,\ (3s+1)\tau\bigr)&(0,0,1)\\
R_s&\bigl(\tau,\ 1-2\tau,\ \tau\bigr)&(0,1,0)\\
O_s&\bigl(4-(3s-1)\tau,\ \tau,\ (3s-2)\tau\bigr)&(3,0,1).
\end{array}
\end{equation}
The $L$ row uses $0\leq s\leq q$; the other rows use $1\leq s\leq q$.
Every coordinate of $u$ is at least $\tau>0$, and $y$ is nonnegative.
The smallest $L$-sector $b$ coordinate is attained at $s=q$ and equals
$\tau$; the smallest $O$-sector $c$ coordinate is attained at $s=1$
and also equals $\tau$. Each diagonal triple receives its value at
both points, so both allocations have total $6q+1$.

We check every triple cover, including all mixed-team triples. If its
$a$ player has an $L$ or $R$ label, its $A$ and $B$ entries are zero.
The possible positive $C$ pairs and the associated two-player sums are
\begin{center}
\begin{tabular}{lll}\toprule
$(b,c)$ labels & $u_b+u_c$ & $y_b+y_c$\\\midrule
$(L_s,L_s)$ & $1-\tau$ & $1$\\
$(R_s,R_s)$ & $1-\tau$ & $1$\\
$(R_s,L_s)$ & $1+(3s-1)\tau$ & $2$\\
$(L_{s-1},O_s)$ & $1-\tau$ & $1$\\\bottomrule
\end{tabular}
\end{center}
All four pair values are one. Adding $u_a=\tau$ or $y_a=0$ covers the
triple. Zero $C$ entries are covered by nonnegativity.

Now suppose its $a$ player is $a_{O_s}$. If at most one of its $A$ and
$B$ entries is positive, the triple value is at most three, covered by
$u_{a_{O_s}}\geq3+4\tau$ and $y_{a_{O_s}}=3$.
If both entries are positive, exactly four cases remain:
\begin{center}
\begin{tabular}{llll}\toprule
$(b,c)$ labels & $w$ & $u_a+u_b+u_c$ & $y_a+y_b+y_c$\\\midrule
$(O_s,O_s)$ & $4$ & $4$ & $4$\\
$(O_s,L_s)$ & $4$ & $4+3\tau$ & $4$\\
$(R_s,O_s)$ & $4$ & $5-3\tau$ & $5$\\
$(R_s,L_s)$ & $5$ & $5$ & $5$\\\bottomrule
\end{tabular}
\end{center}
Since $\tau\leq1/6$, all are covers. The cases also show that every
triple value is at most five.

The matching gives worth at least $6q+1$, while either nonnegative
cover bounds every matching by its total $6q+1$. Hence the true integral
worth is $T=6q+1$ and both $u,y$ belong to the entire core. No fractional
matching value is substituted for coalition worth.

\subsection{A partial matching determines the target}
For each $s=1,\ldots,q$, take the triples
\[
 N_s=(a_{O_s},b_{R_s},c_{L_s}),\qquad
 E_s=(a_{L_{s-1}},b_{L_{s-1}},c_{O_s}).
\]
Their values are five and one, respectively. They form a player-disjoint
partial matching: the used $a$ labels are the $O$ labels and
$L_0,\ldots,L_{q-1}$; the used $b$ labels are the $R$ labels and the same
$L$ labels; the used $c$ labels are $L_1,\ldots,L_q$ and the $O$ labels.
Its total is $6q$. Its unmatched players are exactly
\begin{equation}\label{eq:unmatched}
 U=\{c_{L_0},a_{L_q},b_{L_q}\}
       \cup\{a_{R_s},b_{O_s},c_{R_s}:1\leq s\leq q\},
       \qquad |U|=3q+3.
\end{equation}
Summing the partial-matching inequalities and subtracting from efficiency
gives, at every core point $p$,
\begin{equation}\label{eq:budget}
                       \sum_{a\in U}p_a\leq1.
\end{equation}
Thus a common payoff floor cannot exceed $1/(3q+3)=\tau$. The point
$u$ attains this floor, so the first global leximin level is exactly
$\tau$. At every point of the first optimal face, all coordinates are
at least $\tau$, and $|U|\tau=1$. Inequality~\eqref{eq:budget} therefore
fixes every coordinate of $U$ at $\tau$ on that entire face. In particular,
the full leximin target $c_{L_0}$ is fixed at $\tau$ through every later
refinement. We do not need, and do not assert, that $u$ is the complete
final leximin vector.

The diagonal $L_0$ triple is tight at every core point and has value one.
Nonnegativity implies $H_{c_{L_0}}\leq1$, while the cover $y$ attains
one at that coordinate. Consequently
\[
       H_{c_{L_0}}=1,\qquad x_{c_{L_0}}=\frac{1}{3(q+1)}.
\]
This proves part~(ii) of Theorem~\ref{thm:main}. Ratios are nonnegative,
and choosing arbitrarily large $q$ makes this ratio arbitrarily small;
hence $c_{\rm pair}=0$.

\section{Scope and certificates}
\paragraph{The construction uses the cyclic class essentially.}
The displayed presence of three matrices alone would not rule out another
chain representation of the same tensor. A direct slice obstruction does.
Here the signed chain extension means exactly
$w_{ijk}=\max\{P_{ij}+Q_{jk},0\}$.
If $M_{ij}=\max\{r_i+s_j,0\}$ and both diagonal entries of a $2\times2$
slice are positive, then
$M_{11}+M_{22}\leq M_{12}+M_{21}$: the underlying diagonal sums equal
the underlying off-diagonal sums, and taking positive parts cannot reduce
the latter. For every $q\geq1$, our tensor has the following slices:
\begin{center}
\begin{tabular}{llll}\toprule
Varied sectors & Fixed player & Row; column labels & Slice\\\midrule
$a,b$ & $c_{L_0}$ & $(O_1,L_0);(O_1,L_0)$ & $\left(\begin{smallmatrix}2&1\\0&1\end{smallmatrix}\right)$\\[3pt]
$a,c$ & $b_{L_0}$ & $(O_1,L_0);(O_1,L_0)$ & $\left(\begin{smallmatrix}3&1\\1&1\end{smallmatrix}\right)$\\[3pt]
$b,c$ & $a_{L_0}$ & $(L_0,R_1);(L_0,R_1)$ & $\left(\begin{smallmatrix}1&0\\0&1\end{smallmatrix}\right)$\\\bottomrule
\end{tabular}
\end{center}
Every slice violates the necessary inequality. Whatever sector is chosen
as the middle of a chain, the other two would give a truncated-separable
slice. Thus no sector ordering or within-sector relabeling makes this
tensor a strict chain, even allowing signed adjacent-pair entries with
zero-truncated triple values. This does not exclude a representation
using additional players or a different payoff core.

\paragraph{Reproducible evidence.}
The supplementary certificate package uses exact rational arithmetic.
Its default verifier reconstructs the cyclic family independently and
checks every triple against both covers at $q=1,2,3,4$, along with
matching values, disjointness, the unmatched budget and ideal attainment.
It performs $7208$ triple-cover checks. Frozen input files are also
verified independently against their recorded quantities. An optional
Z3 check at $q=4$ starts only from nonnegativity, efficiency and all raw
triple inequalities and excludes: a larger common floor; a target above
$\tau$ anywhere on the first optimal face; and a target ideal above one.
The optional check also computes full leximin and every ideal from raw
core constraints in three representative chain markets and verifies a
three-payoff-arc cycle certificate for each player.
The chain controls include simplex sharpness and a core with permanent
zeros in which arbitrary first-stage maximizers need not satisfy the
bound. The all-size assertions follow from the proofs above, not from
extrapolating these finite checks. The certificates are exact arithmetic
and solver checks, not a formalization in a proof assistant.

\paragraph{Limits and further questions.}
The construction does not prove a best finite-size ratio or a minimum
number of players. Corollary~\ref{cor:distance} concerns chains and makes
no new claim about opposing selectors in the cyclic class. Classification
for more general interaction graphs, robustness of the construction to
perturbations, and independent historical priority assessment remain
separate questions.

\begin{thebibliography}{99}
\small\setlength{\itemsep}{1pt}
\bibitem{Quint1991} T. Quint. The core of an $m$-sided assignment game.
\emph{Games and Economic Behavior} 3(4):487--503, 1991.
\href{https://doi.org/10.1016/0899-8256(91)90017-9}{doi:10.1016/0899-8256(91)90017-9}.
\bibitem{Lucas1995} W. F. Lucas. Core theory for multiple-sided assignment
games. \emph{Duke Mathematical Journal} 81(1):55--65, 1995.
\href{https://doi.org/10.1215/S0012-7094-95-08106-X}{doi:10.1215/S0012-7094-95-08106-X}.
\bibitem{Stuart1997} H. W. Stuart, Jr. The supplier--firm--buyer game and
its $m$-sided generalization. \emph{Mathematical Social Sciences}
34(1):21--27, 1997.
\href{https://doi.org/10.1016/S0165-4896(96)00830-X}{doi:10.1016/S0165-4896(96)00830-X}.
\bibitem{AtaySolymosi2018} A. Atay and T. Solymosi. On bargaining sets of
supplier-firm-buyer games. \emph{Economics Letters} 167:99--103, 2018.
\href{https://ataatay.github.io/webpage/Atay_Solymosi_2018_EL.pdf}{Author-hosted final article}.
\bibitem{AtayNunez2020} A. Atay and M. N\'u\~nez. Multi-sided assignment
games on $m$-partite graphs. \emph{Annals of Operations Research}
279:271--290, 2019.
\href{https://doi.org/10.1007/s10479-019-03256-5}{doi:10.1007/s10479-019-03256-5}.
\bibitem{NunezRafels2004} M. N\'u\~nez and C. Rafels. Bargained stable
allocations in assignment markets. Barcelona Economics Working Paper 153,
December 2004. \href{https://www.bse.eu/sites/default/files/working_paper_pdfs/153.pdf}{Working version}.
\bibitem{Gangam2025} R. R. Gangam, N. Garg, P. Shahkar, and V. V. Vazirani.
Equitable core imputations for max-flow, MST and $b$-matching games.
arXiv:2403.06037v3, February 2025.
\href{https://arxiv.org/abs/2403.06037v3}{Version inspected}.
\bibitem{Held2003} S. Held, B. Korte, J. Ma\ss berg, M. Ringe, and
J. Vygen. Clock scheduling and clocktree construction for high performance
ASICs. In \emph{Proceedings of ICCAD}, pp.232--239, 2003.
\href{https://www.cs.york.ac.uk/rts/docs/SIGDA-Compendium-1994-2004/papers/2003/iccad03/pdffiles/03d_4.pdf}{Conference text}.
\bibitem{Martinez2020} F. J. Mart\'inez-de-Alb\'eniz, C. Rafels, and N. Ybern.
Assortative multisided assignment games: The extreme core points.
\emph{Games and Economic Behavior} 120:144--153, 2020.
\href{https://doi.org/10.1016/j.geb.2020.01.001}{doi:10.1016/j.geb.2020.01.001}.
\bibitem{Martinez2026} F. J. Mart\'inez-de-Alb\'eniz, C. Rafels, and N. Ybern.
Maximum and minimum core payoffs for supermodular multisided assignment
markets with local monotonicity. \emph{Discrete Applied Mathematics}
379:605--612, 2026.
\href{https://doi.org/10.1016/j.dam.2025.10.006}{doi:10.1016/j.dam.2025.10.006}.
\bibitem{Martinez2023} F. J. Mart\'inez-de-Alb\'eniz, C. Rafels, and N. Ybern.
Extreme core points in multisided assignment problems. Working paper,
SSRN4369524, 2023.
\href{https://doi.org/10.2139/ssrn.4369524}{doi:10.2139/ssrn.4369524}.
\end{thebibliography}
\end{document}
